% ================================================================
%  本章知识地图 · 学习自检 · 知识点—题目对应表
% ================================================================
\usection{本章知识地图与总自检}

\begin{center}
\begin{forest}
  for tree={draw=mainc,fill=lightc,rounded corners=2pt,font=\sffamily\small,align=center,
    grow'=0,l sep=9mm,s sep=2.2mm,inner xsep=4pt,inner ysep=2.5pt,
    edge={draw=auxc,line width=0.6pt},parent anchor=east,child anchor=west,
    edge path={\noexpand\path[\forestoption{edge}] (!u.parent anchor) -- ++(4mm,0) |- (.child anchor)\forestoption{edge label};}}
  [指针\\（第七章）,fill=mainc,text=white,font=\sffamily\bfseries
    [内存与地址\\第一讲
      [按字节编号；变量地址 = 首字节地址]]
    [定义与使用\\第二讲
      [\texttt{int *p = \&a;}\quad \texttt{*p} $\equiv$ \texttt{a}]
      [\texttt{*}：声明中是类型，表达式中是解引用]]
    [指针大小\\第三讲
      [32 位程序 4 字节 / 64 位程序 8 字节，与类型无关]]
    [空指针与野指针\\第四讲
      [空指针：\texttt{nullptr}，可检查，不可解引用]
      [野指针：未初始化 / 乱写地址 / 越界]]
    [const 修饰指针\\第五讲
      [const 在 * 左：锁值]
      [const 在 * 右：锁指向]]
    [指针和数组\\第六讲
      [\texttt{arr} $\to$ \texttt{\&arr[0]}]
      [$p+k$：地址 $+\,k\cdot$\texttt{sizeof(T)}；\texttt{arr[i]} $\equiv$ \texttt{*(arr+i)}]]
    [指针和函数\\第七讲
      [值传递：改副本；地址传递：通过 \texttt{*p} 改实参]]
    [综合：冒泡排序\\第八讲,fill=formc
      [数组传参只传首地址 $\Rightarrow$ 另传长度 $n$]
      [\texttt{i < n-1}，\texttt{j < n-i-1}；比较 $n(n-1)/2$ 次]]
  ]
\end{forest}
\figcap{图 9.1\quad 第七章知识地图：每个分支的叶子是该讲“必须记住的一句话”}
\end{center}

\begin{checkbox}[全章学习自检（不计入正式练习）]
\begin{enumerate}[label=\arabic*.,itemsep=1pt]
  \item 为什么需要指针？说出至少两个“只用变量名做不到”的场景。
  \item \cc{\&}、\cc{*} 分别是什么运算？为什么 \cc{*p} 就是 \cc{a} 本身？
  \item 为什么 \cc{int*} 与 \cc{double*} 大小相同？为什么 \cc{p+1} 走过的字节数却不同？
  \item 空指针和野指针有什么共同点？哪一个可以被程序检查出来？
  \item 看到 \cc{const int * const p}，你能在 3 秒内说出它锁住了什么吗？
  \item \cc{arr[i]} 和 \cc{*(arr+i)} 为什么完全等价？
  \item “地址传递本质上也是值传递”，这句话怎么理解？
  \item 冒泡排序内层循环为什么是 \cc{j < n-i-1}？总共比较几次？
  \item 为什么排序函数需要单独传入数组长度？
  \item 本章知识以后会在哪些课程中再次出现？
\end{enumerate}
\end{checkbox}

\usection{知识点—题目对应表}
下表说明每道正式练习检查的是哪个知识点。做错某题时，可根据“对应讲次”回到正文复习。

\begin{center}
\small
\noindent\begin{tabularx}{\linewidth}{@{}p{0.7cm}p{3.1cm}p{1.3cm}YYY@{}}
\toprule
\rowcolor{lightc}编号 & 核心知识点 & 对应讲次 & A 组基础（1--25） & B 组提升（26--35） & C 组跨学科（36--38） \\
\midrule
K1 & 内存与地址 & 第一讲 & 2，19 & 33 & 38 \\
K2 & 定义、\cc{\&} 与 \cc{*} & 第二讲 & 1，3，11，16，21 & 26，27 & 36，37 \\
K3 & 指针的大小 & 第三讲 & 4，15，17 & 29，33 & — \\
K4 & 空指针 & 第四讲 & 5，12 & — & — \\
K5 & 野指针与越界 & 第四、六讲 & 6，12，24 & 30 & 37 \\
K6 & const 修饰指针 & 第五讲 & 7，8，13，22 & 28，34 & — \\
K7 & 数组与指针算术 & 第六讲 & 9，14，19，24 & 26，28，31，33 & 38 \\
K8 & 值传递与地址传递 & 第七讲 & 10，23 & 27，34 & — \\
K9 & 数组作函数参数 & 第八讲 & 15，18 & 29，34，35 & — \\
K10 & 冒泡排序 & 第八讲 & 20，25 & 30，32，35 & — \\
\bottomrule
\end{tabularx}
\end{center}
每个核心知识点至少被一道基础题覆盖；主要知识点在提升题中再次出现；3 道跨学科题分别借用\textbf{数据结构}（链表）、\textbf{嵌入式系统}（寄存器）、\textbf{图像处理}（像素数组）三个不同背景。
